\section{General Topic Introduction}

\subsection{Project Introduction}

In the context of rapid digital transformation, traditional libraries are gradually transforming into modern online learning resource centers. However, the greatest challenge today extends beyond mere digitization of book titles to managing and exploiting those resources effectively. Existing library management systems often focus primarily on basic administrative tasks, rendering interactions between readers and collections passive and fragmented.

For library managers, manually operating vast resource catalogs—ranging from physical books and journals to digital assets—frequently leads to delays in updating loan statuses and preservation records. The lack of automated analytics tools places heavy pressure on member management and circulation tracking, particularly when forecasting reading demand to optimize collection acquisition.

For readers, searching for materials on legacy systems is often difficult due to rigid keyword matching, lacking flexibility and personalization. Modern users desire a smarter interactive space where they can receive book recommendations tailored to their interests, reading history, and personal behavior without spending excessive time on manual searches. Simultaneously, the demand for immediate assistance via virtual assistants to resolve inquiries regarding rules or document shelf locations is becoming essential.

Recognizing these pressing requirements, the project "Development of an Online Library Management System" was conceived to research, design, and deploy a comprehensive technological platform. This platform not only thoroughly resolves resource management and operational challenges—such as book catalogs, journals, digital assets, and circulation workflows—but also integrates Artificial Intelligence (AI) to enhance collection exploitation. The core focus of the project is creating a smart interactive ecosystem based on a combination of modern administrative models that automate operations and centralize resource management, personalized user experiences through intelligent search tools and automated recommendation systems based on actual reading behavior, and automated support from Chatbots and virtual assistants to deliver immediate query responses and operational guidance to users anytime, anywhere.

By applying software technologies and AI algorithms, the system focuses on optimizing administrative workflows for librarians while providing readers with rapid search capabilities and access to materials tailored to individual needs.

\subsubsection{Motivation for Topic Selection}

During our academic studies and library usage in a university environment, our team observed multiple limitations in existing library systems. Document searches rely heavily on simple keywords, yielding results that often fail to align with user needs. Furthermore, systems lack the capability to recommend materials based on user preferences or borrowing history, requiring readers to spend significant time locating suitable resources.

Additionally, management activities such as tracking loans/returns, updating item statuses, or providing user support remain heavily dependent on manual operations. This not only reduces operational efficiency but also negatively impacts user experience.

Recognizing these realities, our team chose the topic "Development of an Online Library Management System" with the objective of building a modern system capable of automation and intelligent support, contributing to enhanced management efficiency and improved user satisfaction.

\subsection{Project Objectives}

The primary objective of this project is to build an online library management system as a Web Application to effectively support both library administrators and readers in exploiting and utilizing academic resources.

The system is developed to provide a centralized management platform for library operations, including catalog management, loan and return lifecycle tracking, and membership administration. Concurrently, the system aims to enhance user experience through intelligent search tools and recommendation mechanisms based on user preferences and reading history.

Furthermore, the project aims to integrate artificial intelligence technologies to assist users during system interactions. Features such as chatbots or virtual assistants are researched and implemented to deliver instant answers to frequently asked questions, system usage guidance, and rapid document retrieval support.

Through constructing this system, the project strives to create a smart digital library platform that optimizes administrative workflows for library managers while delivering convenient, effective resource discovery and access experiences for readers.

\subsubsection{Overview of Proposed System}

The online library management system is built as a web application, allowing users to access and interact via standard web browsers without requiring additional software installations. Overall, the system comprises core components including the user interface, server-side business logic processing engine, and database storing resource metadata and user profiles.

In addition to basic components, the system integrates artificial intelligence modules to support advanced capabilities such as document recommendations, intelligent search, and chatbot user assistance. These modules operate on data gathered from user behavior and system resource metadata to produce relevant, personalized results.

\subsection{Project Benefits}

The project delivers tangible benefits to library resource management and exploitation.

\subsubsection{Building a Modern Online Library Management System}

The project focuses on developing a web-based library management system enabling resource administration and access through an online environment. The system integrates data processing technologies and artificial intelligence to enhance administrative control and support users during document search and exploitation.

\subsubsection{Effective Support for Library Management}

The system enables management of diverse resource catalogs, including physical books, journals, and digital assets. Features tracking loan-return-preservation statuses, member accounts, and usage analytics empower library staff to operate efficiently. Statistical insights from the system also inform decisions regarding collection acquisition and adjustments.

\subsubsection{Optimizing Reader Resource Exploitation}

Intelligent search capabilities, automated recommendation engines, and virtual assistant features allow users to quickly locate materials matching their learning and research needs. The system also personalizes user experiences based on reading history and search patterns, boosting resource exploitation efficiency.

\subsubsection{Enhancing Interaction Between Users and the Library System}

Integrating chatbots or virtual assistants ensures users receive prompt support when interacting with the system. Inquiries regarding borrowing policies, document location, or system operation are answered immediately, enhancing overall user experience during library interactions.

\subsection{Project Scope and Limitations}

The scope of the project focuses on researching and building an online library management system as a Web Application. Direct system actors comprise: Administrators (Admin), Librarians (Librarian), and Readers (Reader).

\subsubsection{User Groups}

Reader: Primary system users who search, query, and borrow materials. Readers also receive personalized book recommendations, manage loan histories, and interact with the system via search features or chatbots.

Librarian: Responsible for managing catalog entries, processing loan and return transactions, and updating resource statuses. Librarians also utilize administrative tools for data entry and system information control.

Administrator (Admin): Manages the entire system, including parameter configuration, user account management, and overall system monitoring.

The system manages common library resources such as physical books, printed journals, and digital assets in electronic file formats like E-books or PDF documents. Core functionalities encompass catalog management, member account management, loan-return status tracking, document search support, and personalized recommendation generation.

Within the project scope, the system is designed for small to medium-scale libraries, such as departmental libraries, school libraries, or learning resource centers. Target end-users include students, faculty members, and researchers requiring access to academic materials.

AI features are implemented at a foundational level, detailed in the artificial intelligence application scope section below.

\subsubsection{Scope of Artificial Intelligence Application}

Within the project scope, artificial intelligence features are deployed at a foundational level to demonstrate integration potential within a library management system. Specifically, the system utilizes existing models and services to implement functionalities such as behavioral recommendation algorithms, natural language search, and Q\&A chatbots for common inquiries.

Due to time and resource constraints, the project does not focus on building or training AI models from scratch, but primarily leverages existing tools and APIs to deploy and integrate them into the system. This ensures project feasibility while demonstrating a modern technological approach.

\subsubsection{Project Limitations}

Although the system is designed with comprehensive features, certain limitations exist. The system is currently developed exclusively as a web application and does not support native mobile platforms. AI capabilities are implemented at a foundational level and have not reached the optimization of large commercial platforms.

Additionally, the system does not handle extreme large-scale scenarios involving millions of records or deep integration with external third-party library networks. Advanced capabilities will be considered for future research and subsequent software releases.

\subsection{General Topic Significance}

\subsubsection{Practical Significance}

The online library management system developed in this project enhances operational efficiency for libraries in a digital environment. Centralizing resource and user management reduces manual operations while ensuring data accuracy and real-time synchronization.

Furthermore, the system significantly improves reader experiences through intelligent search and recommendation tools. Users can rapidly discover resources tailored to their academic needs without navigating cumbersome search procedures.

Additionally, integrating chatbots or virtual assistants delivers instant user support, reducing the workload for library staff in addressing routine inquiries.

\subsubsection{Scientific Significance}

From a research standpoint, the project clarifies methodologies for building and deploying a library management system using modern software technologies. System design encompasses software architecture engineering, database design, and resource/user management implementation.

The project also explores integrating artificial intelligence models into information systems to personalize user experiences. Applying recommendation algorithms and user-assistive chatbots demonstrates AI's potential in elevating knowledge resource utilization.

By combining web development technologies, data management, and artificial intelligence, the project proposes a smart library management system model that can serve as a benchmark for future research and practical implementations.

In subsequent chapters, the report focuses on system requirement analysis, overall architecture design, and core feature implementation to build a complete library management system.